\titleframe

% Contoh Kode Slide
%\section{Introduction}
%\begin{frame}{\insertsectionhead}
%	\framesubtitle{A short introduction to Trigon}
%	\themename is a modern, elegant and versatile theme for Beamer, inspired by
%	the
%	\href{https://github.com/matze/mtheme}{\textsc{metropolis} theme} from Matthias
%	Vogelgesang.
%	\vfill
%	\themename comes with lots of nice extra features
%	\begin{itemize}
%		\item Multiple style variations for title, section and normal slides
%		\item Simple customization of theme colors
%		\item Lots of convenient options to tweak the design
%	\end{itemize}
%\end{frame}

% Contoh Kode Tabel
%  \begin{table}[H]
% \centering
% \caption{A nice table example}
% \begin{tabular}{@{} lccc @{}}
%    \toprule
%    & \textbf{Velocity} & \textbf{Angle}  & \textbf{Vertical force} \\
%    & $U$ & $\alpha$  & $F_z$ \\
%    & [m/s] & [$^\circ$]  & [N] \\
%    \midrule
%    2D simulation  & 9 & 2 & 9.23 \\
%    3D simulation  & 10.0 & 3 & 15.039 \\
%    Experiment A   & 11.31 & 2.5 & 13.2 \\
%    Experiment B   & 11.26 & 2.7 & 12.6 \\
%    Experiment C   & 11.33 & 2.47 & 13.6 \\
%    \bottomrule
%  \end{tabular}
%\end{table}

% Contoh Kode Gambar
%  \begin{figure}[ht!]
%	\begin{subfigure}[b]{0.3\textwidth}
%		\frame{\includegraphics[width=\textwidth]{layout_example-03.jpg}}
%		\caption*{plain}
%	\end{subfigure}
%	\hspace{\fill}
%	\begin{subfigure}[b]{0.3\textwidth}
%		\frame{\includegraphics[width=\textwidth]{layout_example-02.jpg}}
%		\caption*{style1}
%	\end{subfigure}
%	\hspace{\fill}
%	\begin{subfigure}[b]{0.3\textwidth}
%		\frame{\includegraphics[width=\textwidth]{layout_example-01.jpg}}
%		\caption*{style2 (default)}
%	\end{subfigure}
%\end{figure}

% Contoh Kode Blocks
% \begin{block}{Regular block}
%	Just a regular block
%\end{block}
%\begin{alertblock}{Alert block}
%	Some important thing
%\end{alertblock}
%\begin{exampleblock}{Example block}
%	No difference with regular block to avoid excessive distraction
%\end{exampleblock}

\setbeamertemplate{frame footer}{Alauddin Maulana Hirzan, M. Kom - Universitas Semarang}

\section{Teknologi Konektivitas \textit{Internet of Things}}
\begin{frame}{\insertsectionhead}
	\framesubtitle{Konektivitas Perangkat \#1}
	\justifying
	Teknologi \textit{Internet of Things} merupakan teknologi yang sangat bergantung kepada konektivitas perangkat. Tanpanya perangkat tidak akan bisa berkomunikasi dengan perangkat lainnya atau dengan pengguna akhir.
	\vfill
	Sehingga untuk bisa memanfaatkan teknologi \textit{Internet of Things} secara maksimal, teknologi konektivitas \textit{Internet of Things} harus diperhatikan. Perangkat-perangkat tertentu ada yang sudah dilengkapi dengan teknologi jaringan, dan sebagian ada yang belum.
\end{frame}

\begin{frame}{\insertsectionhead}
	\framesubtitle{Konektivitas Perangkat \#2}
	\justifying
	Di pasaran \textit{Internet of Things} banyak teknologi yang bisa digunakan untuk komunikasi pera-ngkat, berupa:
	\begin{multicols}{2}
		\begin{itemize}
			\item \textit{Bluetooth Low Energy (BLE)} 
			\item \textit{Wi-Fi}
			\item \textit{Z-Wave}
			\item \textit{Thread}
			\item \textit{ZigBee}
			\item \textit{Satellite}
			\item \textit{SigFox}
			\item \textit{Low Power Wide Area (LPWA)}
			\item \textit{Long Range Wide Area Network (LoRaWAN)}
		\end{itemize}
	\end{multicols}
\end{frame}

\begin{frame}{\insertsectionhead}
	\framesubtitle{Konektivitas Perangkat \#2}
	\justifying
	Perangkat seperti \textbf{BLE} memiliki jarak yang dekat, \textbf{Wi-Fi} memiliki jangkauan yang sedang, dan teknologi seperti \textbf{LoRa} memiki radius yang sangat jauh. Sehingga untuk pemanfaatannya tergantung dari situasi dan lokasi.
	\vfill
	\begin{block}{Info}
		\textbf{\textit{Long Range Wide Area Network}} atau disingkat sebagai \textbf{\textit{LoRa}} merupakan teknologi komunikasi jarak jauh yang dapat menjangkau daerah tanpa komunikasi efektif seperti gunung.
	\end{block}
\end{frame}

\begin{frame}{\insertsectionhead}
	\framesubtitle{Konektivitas Perangkat \#2}
	\justifying
	  \begin{figure}[ht!]
			\begin{subfigure}[b]{0.5\textwidth}
					\frame{\includegraphics[width=\textwidth]{bluetooth.png}}
					\caption*{Bluetooth}
				\end{subfigure}
			\hspace{\fill}
			\begin{subfigure}[b]{0.19\textwidth}
					\frame{\includegraphics[width=\textwidth]{wifi.png}}
					\caption*{Wireless Fidelity}
				\end{subfigure}
			\hspace{\fill}
			\begin{subfigure}[b]{0.25\textwidth}
					\frame{
\includegraphics[width=\textwidth]{lora.png}}
					\caption*{Long Range}
				\end{subfigure}
		\end{figure}
\end{frame}

\begin{frame}{\insertsectionhead}
	\framesubtitle{Teknologi \textit{Bluetooth}}
	\justifying
	Dilansir dari \textbf{bluetooh.com}:
	\vfill
	\begin{quote}
		\begin{itemize}
			\item Bluetooth Klasik adalah radio berdaya rendah yang mengalirkan data melalui 79 saluran dalam pita frekuensi industri, ilmiah, dan medis (ISM) sebesar 2,4 GHz tanpa lisensi.
			\item Radio Bluetooth Low Energy (LE) dirancang untuk pengoperasian dengan daya yang sangat rendah. Mengirimkan data lebih dari 40 saluran dalam pita frekuensi ISM 2,4 GHz tanpa izin.
		\end{itemize}
	\end{quote}
\end{frame}

\begin{frame}{\insertsectionhead}
	\framesubtitle{Lisensi?}
	\justifying
	Dalam penggunaan gelombang radio membutuhkan perizinan dari pemerintah di negara manapun. Karena berkaitan dengan ruang publik agar tidak dimonopoli oleh satu pihak. Namun kasus khusus \textit{Bluetooth} yang memiliki jangkauan yang rendah <10m, sehingga tidak memerlukan lisensi atau perizinan dari pemerintah.
	\vfill
	\begin{block}{Info}
		Alokasi spektrum gelombang radio di Indonesia sudah diatur dalam Peraturan Menteri nomor 29 tahun 2009 oleh Kementerian Komunikasi dan Informatika Indonesia yang juga mengacu ke \textit{International Telecommunication Union}
	\end{block}
\end{frame}

\begin{frame}{\insertsectionhead}
	\framesubtitle{Spesifikasi \textit{BLE}}
	\justifying
	\begin{multicols}{2}
		\begin{itemize}
			\item Pita ISM 2,4 GHz
			\item 3 saluran iklan
			\item 37 saluran data
			\item Point-to-Point, Broadcast, Mesh
			\item Modulasi GFSK
			\item Direction: Direction Finding (AoA/AoD)
		\end{itemize}
	\end{multicols}
	\vfill
	 \begin{figure}[ht!]
	\begin{subfigure}[b]{0.65\textwidth}
		\frame{\includegraphics[width=\textwidth]{survey.png}}
	\end{subfigure}
\end{figure}
\end{frame}

\begin{frame}{\insertsectionhead}
	\framesubtitle{Implementasi Teknologi Bluetooth}
	\justifying
	Teknologi Bluetooth dapat dimanfaatkan untuk berbagai macam hal:
	\begin{itemize}
		\item \textbf{Audio Streaming}
		\item \textbf{Transfer data}h banyak lagi.
		\item \textbf{Perangkat Jaringan}
		\item \textbf{Layanan Lokasi}
	\end{itemize}
\end{frame}

\begin{frame}{\insertsectionhead}
	\framesubtitle{\textit{BLE Beacon Tracking}}
	\justifying
	 \begin{figure}[ht!]
		\begin{subfigure}[b]{0.7\textwidth}
			\frame{\includegraphics[width=\textwidth]{ble-floorplan.jpg}}
		\end{subfigure}
	\end{figure}
\end{frame}

\begin{frame}{\insertsectionhead}
	\framesubtitle{Teknologi \textit{Wi-Fi}}
	\justifying
	Wi-Fi (juga ditulis Wifi atau WiFi) adalah sebuah teknologi yang memanfaatkan peralatan elektronik untuk bertukar data secara nirkabel (menggunakan gelombang radio) melalui sebuah jaringan komputer, termasuk koneksi Internet berkecepatan tinggi.
	\vfill
	\begin{block}{Info}
		Wi-Fi =//= Wireless Fidelity	
	\end{block}
\end{frame}

\begin{frame}{\insertsectionhead}
	\framesubtitle{Miskonsepsi \textit{Wi-Fi}}
	\justifying
	Wi-Fi bukan \textit{Wireless Fidelity}, miskonsepsi yang masih dibawa hingga sekarang. Istilah Wi-Fi diciptakan pada tahun 1999 oleh Interbrand, perusahaan konsultan merek yang disewa oleh Wi-Fi Alliance, organisasi yang bertanggung jawab untuk promosi teknologi Wi-Fi dan sertifikasi produk Wi-Fi.
	\vfill
	Kesalahpahaman Wi-Fi adalah singkatan dari "Wireless Fidelity", sebenarnya tidak sengaja dibuat oleh Wi-Fi Alliance sendiri ketika Wi-Fi diperkenalkan. 
\end{frame}

\begin{frame}{\insertsectionhead}
	\framesubtitle{Teknologi Wi-Fi\textsuperscript{\textcopyright} \#1}
	\justifying
	Wi-Fi\textsuperscript{\textcopyright} adalah
	\begin{itemize}
		\item Teknologi komunikasi nirkabel yang paling umum digunakan
		\item Media utama untuk lalu lintas internet global
		\item Penggerak nilai ekonomi global \$3,3 triliun USD
		\item Berkembang, dengan lebih dari 4 miliar perangkat dikirimkan setiap tahun dan 16 miliar perangkat digunakan*
	\end{itemize}	
\end{frame}

\begin{frame}{\insertsectionhead}
	\framesubtitle{Teknologi Wi-Fi\textsuperscript{\textcopyright} \#2}
	\justifying
	Sama seperti teknologi komunikasi gawai, teknologi nirkabel ini juga memiliki generasi hingga 6G. Sedangkan Wi-Fi generasi 7 masih dalam pengembangan.
	\vfill
	Wi-Fi CERTIFIED 6™, program sertifikasi industri berdasarkan standar IEEE 802.11ax, menyediakan kapasitas, efisiensi, jangkauan, dan kinerja yang dibutuhkan oleh pengguna saat ini di lingkungan Wi-Fi® yang paling menuntut.
	\vfill
	\begin{block}{Info}
		Raspberry Pi versi 4 menggunakan teknologi Wi-Fi generasi 5 dengan standar 802.11ac.
	\end{block}
\end{frame}

\begin{frame}{\insertsectionhead}
	\framesubtitle{\textit{Wi-Fi 6E}}
	\justifying
	\begin{figure}[ht!]
		\begin{subfigure}[b]{0.7\textwidth}
			\frame{\includegraphics[width=\textwidth]{wifi6e.png}}
		\end{subfigure}
	\end{figure}
\end{frame}

\begin{frame}{\insertsectionhead}
	\framesubtitle{Teknologi \textit{Z-Wave}}
	\justifying
	Dilansir dari \textbf{z-wave.com}:
	\vfill
	\begin{quote}
		Z-Wave didasarkan pada topologi jaringan mesh. Ini berarti setiap perangkat (non-baterai) yang terpasang di jaringan menjadi pengulang sinyal. Sehingga, semakin banyak perangkat yang Anda miliki di rumah, semakin kuat jaringannya.
		\vfill
		Sementara Z-Wave memiliki jangkauan 100 meter atau 328 kaki di udara terbuka, bahan bangunan mengurangi jarak tersebut, disarankan untuk memiliki perangkat Z-Wave kira-kira setiap 30 kaki, atau lebih dekat untuk efisiensi maksimum.
	\end{quote}
\end{frame}

\begin{frame}{\insertsectionhead}
	\framesubtitle{Teknologi \textit{ZigBee}}
	\justifying
	Sedangkan untuk \textbf{ZigBee}:
	\vfill
	\begin{quote}
		Zigbee adalah spesifikasi berbasis IEEE 802.15.4 untuk rangkaian protokol komunikasi tingkat tinggi yang digunakan untuk membuat jaringan area pribadi dengan radio digital kecil berdaya rendah, seperti untuk otomatisasi rumah, pengumpulan data perangkat medis, dan daya rendah lainnya. kebutuhan bandwidth rendah, dirancang untuk proyek skala kecil yang membutuhkan koneksi nirkabel
	\end{quote}
\end{frame}

\begin{frame}{\insertsectionhead}
	\framesubtitle{Teknologi \textit{ZigBee}}
	\justifying
	\begin{figure}[ht!]
		\begin{subfigure}[b]{0.4\textwidth}
			\frame{\includegraphics[width=\textwidth]{zuno.jpg}}
			\caption*{Z-Uno}
		\end{subfigure}
	\hspace{50pt}
	\begin{subfigure}[b]{0.4\textwidth}
		\frame{\includegraphics[width=\textwidth]{zigbee.jpg}}
		\caption*{ZigBee}
	\end{subfigure}
	\end{figure}	
\end{frame}

\begin{frame}{\insertsectionhead}
	\framesubtitle{Teknologi \textit{Low-Power Wide-Area (LPWA) }}
	\justifying
	LPWA adalah singkatan dari Low-Power Wide-Area. Ini tidak mengacu pada satu teknologi tertentu, melainkan berfungsi sebagai istilah umum untuk jaringan apa pun yang dirancang untuk berkomunikasi secara nirkabel dengan daya yang lebih rendah daripada jaringan lain seperti seluler, satelit, atau WiFi.
	\vfill
	\begin{block}{Info}
		Jaringan seluler modern memiliki kecapatan hingga gigabit per detik dengan LTE Advanced dan jaringan 5G yang akan datang. Sedangkan, jaringan LPWA mentransmisikan jauh lebih sedikit—seringkali hanya beberapa kilobit per saluran.
	\end{block}
\end{frame}

\begin{frame}{\insertsectionhead}
	\framesubtitle{Implementasi \textit{Low-Power Wide-Area (LPWA) }}
	\justifying
	Bandwidth jaringan LPWA yang sangat terbatas tidak cocok untuk:
	\vfill
	\begin{itemize}
		\item Suara
		\item Video
		\item Audio
		\item Teks
	\end{itemize}
	 \vfill
	 Dengan demikian, jaringan LPWA hanya digunakan untuk perangkat Internet of Things (IoT), dan komunikasi M2M (mesin-ke-mesin).
\end{frame}

\begin{frame}{\insertsectionhead}
	\framesubtitle{Produk \textit{Low-Power Wide-Area (LPWA) }}
	\justifying
	Sigfox, jaringan LPWA eksklusif yang ditawarkan oleh perusahaan Prancis. Produk ini menggunakan frekuensi tanpa izin di pita 868 MHz atau 902 MHz. Dicirikan sebagai transmisi radio ultra-narrowband, ia menawarkan jangkauan jarak jauh, tetapi dengan kecepatan transfer data yang rendah.
	\vfill
	\begin{figure}[ht!]
	\begin{subfigure}[b]{0.6\textwidth}
		\frame{\includegraphics[width=\textwidth]{sigfox.png}}
	\end{subfigure}
\end{figure}	
\end{frame}

\begin{frame}{\insertsectionhead}
	\framesubtitle{Penggunaan \textit{Low-Power Wide-Area (LPWA) }}
	\justifying
	LPWA bisa digunakan untuk fungsi-fungsi tertentu seperti:
	\vfill
	\begin{itemize}
		\item Manajemen parkir
		\item Meter air dan saluran pipa
		\item Pengiriman cerdas
		\item Penerangan jalan dan penerangan jalan raya
	\end{itemize}
\end{frame}

\begin{frame}{\insertsectionhead}
	\framesubtitle{Teknologi \textit{Long Range (LoRa) }}
	\justifying
	LoRa merupakan salah satu bentuk produk dari LPWA yang sudah disertifikasi oleh LoRa Alliance. Sehingga dalam penggunaannya bisa berbeda dari produk LPWA lainnya.
	\vfill
	Dilansir dari \textbf{lora-alliance.org}:
	\vfill
	\begin{quote}
		LoRaWAN® adalah protokol jaringan Low Power Wide Area (LPWA) yang dirancang untuk menghubungkan 'sesuatu' yang dioperasikan dengan baterai secara nirkabel ke internet di jaringan regional, nasional, atau global.
	\end{quote}
\end{frame}

\begin{frame}{\insertsectionhead}
	\framesubtitle{Arsitektur \textit{LoRaWAN}}
	\justifying
	\begin{figure}[ht!]
		\begin{subfigure}[b]{0.85\textwidth}
			\frame{\includegraphics[width=\textwidth]{lorawan.png}}
		\end{subfigure}
	\end{figure}	
\end{frame}

\begin{frame}{\insertsectionhead}
	\framesubtitle{Perangkat \textit{LoRaWAN}}
	\justifying
	\begin{figure}[ht!]
		\begin{subfigure}[b]{0.65\textwidth}
			\frame{\includegraphics[width=\textwidth]{loraesp.jpg}}
		\end{subfigure}
	\end{figure}	
\end{frame}

\setbeamertemplate{background} 
{
	\includegraphics[width=\paperwidth,height=\paperheight]{thank-you.jpg}
}
\begin{frame}[plain]
\end{frame}